\capitulo{5}{Aspectos relevantes del desarrollo del proyecto}

\section{Planteamiento y organización del desarrollo}

El desarrollo de este trabajo no comenzó desde una aplicación vacía, sino a partir de un editor de diagramas Entidad-Relación ya existente. Antes de añadir nuevas funcionalidades fue necesario comprender cómo estaba construido el sistema, qué partes podían aprovecharse y qué aspectos convenía revisar para poder seguir ampliándolo con seguridad.

Una de las ideas principales del proyecto ha sido mantener la coherencia entre las distintas partes del editor. No basta con que un elemento se dibuje correctamente en el lienzo: también debe quedar representado en el modelo interno, ya que esa información se utiliza después para guardar el diagrama, reconstruirlo, validarlo, transformarlo al modelo relacional y generar SQL.

Por ello, el desarrollo se planteó de forma incremental, utilizando ramas, commits, issues y milestones de GitHub para separar tareas y mantener trazabilidad. Esta organización resultó especialmente útil al trabajar sobre una aplicación heredada, donde un cambio podía afectar simultáneamente a la interfaz, el modelo interno, la persistencia, la validación o la generación de SQL.

La Tabla~\ref{tab:comparativa_tfg_anterior_actual} resume la diferencia entre el punto de partida heredado y la aportación realizada en este trabajo. La comparación no pretende desvalorizar la aplicación previa, sino delimitar qué partes se han reutilizado y qué aspectos se han ampliado, estabilizado o documentado durante este TFG.

\begin{table}[htbp]
    \centering
    \small
    \caption{Comparación entre la versión heredada y la versión desarrollada en este TFG.}
    \label{tab:comparativa_tfg_anterior_actual}
    \begin{tabular}{p{0.24\textwidth} p{0.28\textwidth} p{0.36\textwidth}}
        \toprule
        \textbf{Aspecto} & \textbf{Versión heredada} & \textbf{Aportación de este TFG} \\
        \midrule
        Interfaz del editor & Controles básicos sobre el lienzo. & Reorganización del panel lateral, branding UBU, selector de idioma, acciones contextuales, textos de apoyo e información del proyecto. \\
        Modelo interno & Representación inicial de entidades, atributos y relaciones. & Ampliación del metamodelo para entidades débiles, atributos compuestos y multivaluados, ternarias, roles, relaciones identificadoras e ISA inicial. \\
        Validación & Comprobaciones básicas antes de generar resultados. & Reglas ampliadas, mensajes agrupados y localización más precisa del elemento relacionado con cada error cuando es posible. \\
        Transformación relacional & Transformación básica hacia tablas y claves. & Revisión de casos complejos: multivaluados, entidades débiles, ternarias, roles, claves compuestas e ISA con clave heredada. \\
        Generación SQL & Generación inicial de sentencias SQL. & SQL más legible, ordenación de tablas para reducir \texttt{ALTER TABLE}, simplificación de claves foráneas, restricciones \texttt{UNIQUE} y normalización de identificadores. \\
        Persistencia y usabilidad & Persistencia e interacción básicas. & Importación/exportación JSON, reconstrucción visual, selección múltiple, ajuste de vista, exportación visual y estructuras de ejemplo. \\
        Pruebas y documentación & Base inicial de desarrollo y documentación previa. & Ampliación de pruebas, revisión manual de casos complejos y documentación técnica, de usuario y de diseño actualizada. \\
        \bottomrule
    \end{tabular}
\end{table}

\section{Análisis y estabilización del sistema heredado}

Antes de incorporar ampliaciones, se realizó una revisión del sistema heredado. La aplicación ya permitía trabajar con diagramas E/R básicos, pero era necesario comprender la relación entre el lienzo de mxGraph, el modelo interno y las operaciones posteriores de validación y generación.

Durante este análisis se identificaron varios puntos críticos. En primer lugar, la aplicación funcionaba en dos niveles: la parte visual, gestionada mediante mxGraph, y el modelo interno, utilizado por la lógica de la aplicación. Si ambos niveles no se actualizan de forma coherente, puede ocurrir que el usuario vea un diagrama en pantalla que no coincida con la información que después se valida o se transforma a SQL.

También se revisaron los mecanismos de carga, guardado, importación y exportación. En una aplicación de este tipo, la persistencia no consiste únicamente en almacenar un dibujo, sino en conservar suficiente información semántica para reconstruir el diagrama y seguir trabajando con él. Por ello, se normalizaron estructuras internas, se revisó la reconstrucción visual y se mejoró el tratamiento de identificadores al importar o combinar diagramas.

Ejemplos concretos de estabilización fueron la revisión de la sincronización entre celdas visuales y modelo interno, la normalización de diagramas con campos incompletos, el remapeo de identificadores al importar JSON y la mejora de los mensajes de validación para indicar, cuando es posible, el elemento relacionado con cada problema.

\section{Evolución del modelo interno y persistencia}

El modelo interno del diagrama actúa como estructura central del editor. A partir de él se validan los elementos, se transforma el diagrama al modelo relacional, se genera SQL y se reconstruye el contenido al cargar un fichero o recuperar el estado guardado.

La Figura~\ref{fig:modelo_interno_editor} resume esta relación entre el lienzo visual, el modelo interno y las operaciones principales de la aplicación.

\imagen{modelo_interno_editor}{Relación entre el lienzo del editor, el modelo interno y las operaciones principales de la aplicación.}{0.80}

Durante el proyecto se amplió este modelo para incorporar información que antes no era necesaria o no estaba representada de forma suficiente. Esto incluye propiedades asociadas a entidades débiles, discriminantes, atributos compuestos y multivaluados, relaciones ternarias, roles, relaciones identificadoras y jerarquías ISA.

También se revisó la persistencia local y la importación/exportación JSON. La exportación permite guardar un diagrama en un formato reutilizable, mientras que la importación permite recuperar diagramas previos. Además, se incorporaron modos de importación para reemplazar el diagrama actual o combinarlo con el existente, evitando conflictos mediante el remapeo de identificadores y el desplazamiento visual del bloque importado.

\section{Ampliación de los elementos Entidad-Relación soportados}

Una parte relevante del trabajo consistió en ampliar los elementos del modelo Entidad-Relación soportados por la aplicación. Estas ampliaciones no se limitaron a añadir nuevas formas visuales, sino que exigieron adaptar el modelo interno, la validación, la persistencia, la transformación relacional y la generación SQL.

Entre los elementos incorporados o revisados se encuentran:

\begin{itemize}
    \item \textbf{Entidades débiles y relaciones identificadoras}: se añadió soporte para entidades que dependen de una entidad fuerte y utilizan un discriminante como parte de su identificación.

    \item \textbf{Atributos compuestos y multivaluados}: se revisó su representación visual e interna, así como su transformación posterior al modelo relacional.

    \item \textbf{Relaciones ternarias y roles}: se amplió el soporte de relaciones de grado superior, incluyendo participantes repetidos mediante roles para evitar ambigüedades.

    \item \textbf{Jerarquías ISA}: se incorporó un soporte inicial para generalizaciones y especializaciones, con herencia de clave hacia las tablas especializadas.
\end{itemize}

En el caso de ISA, el alcance se mantuvo deliberadamente limitado. La aplicación permite una generalización y una o varias especializaciones, pero no incorpora restricciones de totalidad o parcialidad, solapamiento o exclusividad, discriminantes ni herencia múltiple. Esta decisión permitió ofrecer una primera base funcional sin aumentar excesivamente la complejidad del proyecto.

\section{Validación, transformación relacional y generación de SQL}

La validación del diagrama es una de las partes más importantes de la herramienta, ya que evita generar resultados a partir de modelos incompletos o incoherentes. Entre las comprobaciones realizadas se incluyen reglas generales, como evitar elementos sin nombre o nombres repetidos cuando pueden producir ambigüedad, y reglas específicas del modelo E/R.

Por ejemplo, se revisa que cada entidad tenga atributos y algún mecanismo de identificación válido: una clave primaria en entidades fuertes, un discriminante en entidades débiles o una clave heredada cuando la entidad actúa como especialización ISA. También se comprueba que las relaciones tengan participantes y cardinalidades configurados de forma válida.

La Figura~\ref{fig:dialogo_validacion} muestra un ejemplo de errores detectados durante la validación. En la versión heredada, la validación podía indicar el tipo de problema, pero no siempre permitía localizar claramente dónde se producía. En la versión actual, cuando la información está disponible, el mensaje añade el elemento concreto relacionado con el error, por ejemplo indicando qué entidad no tiene clave primaria.

\imagen{dialogo_validacion}{Diálogo de validación con errores agrupados del diagrama.}{0.75}

Tras la validación, la aplicación transforma el diagrama al modelo relacional. Esta transformación incluye la creación de tablas para entidades, relaciones muchos a muchos, atributos multivaluados, relaciones ternarias, entidades débiles y especializaciones ISA. También se revisa la propagación de claves primarias y claves foráneas según las cardinalidades y el tipo de elemento representado.

La generación de SQL se apoya en esa representación relacional. El SQL generado utiliza una sintaxis sencilla, basada en construcciones habituales del estándar ANSI/ISO SQL y compatible con PostgreSQL~\cite{postgresqlCreateTable}. En particular, se emplean sentencias \texttt{CREATE TABLE}, claves primarias, claves foráneas, restricciones \texttt{UNIQUE} y restricciones \texttt{NOT NULL}.

Durante la revisión del generador se buscó producir un script más legible. Para ello, las claves foráneas se incluyen dentro de las sentencias \texttt{CREATE TABLE} cuando las dependencias entre tablas pueden resolverse mediante el orden de creación. Las sentencias \texttt{ALTER TABLE} se mantienen como respaldo para dependencias circulares que no pueden resolverse solo mediante la ordenación.

También se simplificó la representación de las claves foráneas que referencian la clave primaria completa de la tabla destino. Por ejemplo, en lugar de generar \texttt{REFERENCES Oficinas(idOficina)}, puede generarse \texttt{REFERENCES Oficinas}. Además, se normalizan identificadores para evitar problemas con espacios, acentos u otros caracteres poco adecuados para SQL.

La Figura~\ref{fig:sql_generado} muestra un ejemplo de script SQL generado por la aplicación.

\imagen{sql_generado}{Ejemplo de script SQL generado a partir de un diagrama validado.}{0.80}

El resultado debe entenderse como una traducción académica y simplificada del modelo diseñado. Por ejemplo, los atributos se traducen con un tipo general \texttt{VARCHAR(40)}, suficiente para representar la estructura del esquema, pero sin entrar en una configuración detallada de tipos físicos.

\section{Refinamientos de usabilidad del editor}

Además de las ampliaciones del modelo, se incorporaron mejoras orientadas a facilitar el uso del editor. Entre ellas se encuentran la reorganización del panel lateral, el branding UBU, el selector de idioma, la selección múltiple visible, el borrado y movimiento conjunto de elementos, el ajuste de vista, la exportación visual, la importación JSON y la generación de estructuras de ejemplo.

También se revisaron mensajes de confirmación, validación y error para que resultaran más claros. Estas mejoras no modifican la teoría del modelo Entidad-Relación, pero reducen fricción durante la edición y ayudan a que la herramienta resulte más usable en un contexto docente.

\section{Pruebas y revisión manual}

La comprobación del proyecto se realizó combinando pruebas automáticas y revisión manual. Las pruebas unitarias se utilizaron para validar funciones de dominio, transformación y generación, mientras que las pruebas end-to-end permitieron comprobar flujos completos de uso sobre la aplicación en navegador.

Además, se realizó revisión manual en aquellos casos donde era importante comprobar el comportamiento visual o la interacción del usuario, como la reconstrucción del lienzo, la selección múltiple, la importación de diagramas, la edición de elementos o la generación visual de estructuras.

En las comprobaciones recientes se ejecutaron los siguientes comandos:

\begin{verbatim}
npm run lint
npm test
npm run test:e2e
npm run build
\end{verbatim}

Estas comprobaciones permitieron verificar que el proyecto mantenía un estado estable antes de la fase final de documentación y revisión.

\section{Dificultades y decisiones técnicas relevantes}

El desarrollo presentó varias dificultades derivadas del carácter heredado y visual de la aplicación. Una de las principales fue mantener la sincronización entre mxGraph y el modelo interno, ya que el lienzo representa elementos gráficos, mientras que la aplicación necesita conservar información semántica para validar, transformar y generar SQL.

Otra dificultad fue decidir el alcance de las ampliaciones. Algunos elementos, como las jerarquías ISA, podían crecer mucho en complejidad si se incorporaban todas sus variantes teóricas. Por ello, se optó por un soporte inicial controlado, dejando fuera restricciones de totalidad o parcialidad, solapamiento o exclusividad, discriminantes y herencia múltiple.

También fue necesario equilibrar la legibilidad del SQL generado con la complejidad de las dependencias entre tablas. Por esta razón se priorizó incluir claves foráneas dentro de \texttt{CREATE TABLE} cuando el orden de creación lo permite, manteniendo \texttt{ALTER TABLE} únicamente como mecanismo de respaldo para dependencias circulares.

En conjunto, estas decisiones reflejan el enfoque seguido durante el proyecto: evolucionar la aplicación de forma mantenible, ampliando su utilidad docente sin convertirla en una herramienta profesional completa de diseño de bases de datos.